% !PRINTED 144
\weberpage{144}
%<<CONT
and consequently
\weq{1}{j \geqq \frac{\Pi(n)}{[\Pi(s)]^r\,\Pi(r)}.}
It is easy to see that this number, when neither of the
factors $r, s$ is equal to 1, is greater than $n$. This follows
if we write the quotient thus:
\[
\frac{\Pi(n)}{[\Pi(s)]^r\,\Pi(r)} = \frac{(r+1)(r+2)\dots n}{(2.3\dots s)^r} =
\]
\[
\frac{n}{2}\cdot\left(\frac{r+1}{2}\cdot\frac{r+2}{2}\dots\frac{2r-1}{2}\right)\left(\frac{2r}{3}\cdot\frac{2r+1}{3}\dots\frac{3r-1}{3}\right)\dots
\]
\[
\left(\frac{(s-1)r}{s}\cdot\frac{(s-1)r+1}{s}\dots\frac{n-1}{s}\right).
\]
For according to this
\weq{2}{j > \frac{n}{2}\left(\frac{r+1}{2}\right)^{r-1}\left(\frac{2r}{3}\right)^r \dots \left(\frac{(s-1)r}{s}\right)^r.}
The factors $\dfrac{r+1}{2}, \dfrac{2r}{3}, \dots \dfrac{(s-1)r}{s}$ are all equal to or
greater than 1, for we have
\[
\frac{(h-1)r}{h} - 1 = \frac{(h-1)(r-1)-1}{h},
\]
which is therefore always positive, and the first factor
\[
\left(\frac{r+1}{2}\right)^{r-1}
\]
is greater than 2 when $r > 2$. But if $r = 2$, then the
product of the two factors is:
\[
\left(\frac{r+1}{2}\right)^{r-1}\left(\frac{2r}{3}\right)^r = \frac{3}{2}\left(\frac{4}{3}\right)^2 = \frac{8}{3} > 2,
\]
and consequently in all circumstances $j > n$.

If, secondly, the group $Q$ is intransitive, and the
system of the $n$ figures splits into two systems of $a$ and $b$ figures respectively, so
that the figures of these two systems are permuted by $Q$ only among
 themselves, and $n = a + b$, then all the permutations
 of $Q$ are contained in the group that consists of all the permutations
 of the $a$ and of the $b$ figures; that is, the degree
of $Q$ is equal to or less than $\Pi(a)\Pi(b)$, and consequently the
index $j$ of $Q$
\weq{3}{j \geqq \frac{\Pi(n)}{\Pi(a)\,\Pi(b)} = \frac{n}{1}\,\frac{n-1}{2}\dots\frac{b+1}{a}.}
